\documentclass{article}

% packages
\usepackage{amsmath, amsthm, thmtools, amsfonts, amssymb, luacode, catchfile, tikzducks, hyperref, ifthen}
\ifcsname c@kobocompile\endcsname
	\usepackage[a5paper, total={1072pt, 1448pt}, margin=10pt, includeheadfoot]{geometry} % set page margins
\else
	\usepackage[a4paper, margin=50pt, includeheadfoot]{geometry}
\fi
\usepackage[shortlabels]{enumitem}
\usepackage[skip=3pt, indent=0pt]{parskip}

% language
\usepackage[bidi=basic, layout=tabular, provide=*]{babel}
\ifcsname c@english\endcsname
	\babelprovide[main, import]{english}
\else
	\babelprovide[main, import]{hebrew}
	\babelprovide{rl}
\fi
%\babelfont{rm}{Libertinus Serif}
\babelfont{rm}[Renderer=Harfbuzz]{Libertinus Serif}
\babelfont{sf}{Libertinus Sans}
\babelfont{tt}{Libertinus Mono}

% style
\AddToHook{cmd/section/before}{\clearpage}	% Add line break before section
\linespread{1.3}
\setcounter{secnumdepth}{0}		% Remove default number tags from sections, this won't do well with theorems
\AtBeginDocument{\setlength{\belowdisplayskip}{3pt}}
\AtBeginDocument{\setlength{\abovedisplayskip}{3pt}}
\graphicspath{ {../images/} }

% operators
\DeclareMathOperator\cis{cis}
\DeclareMathOperator\Sp{Sp}
\DeclareMathOperator\tr{tr}
\DeclareMathOperator\im{Im}
\DeclareMathOperator\re{Re}
\DeclareMathOperator\diag{diag}
\DeclareMathOperator*\lowlim{\underline{lim}}
\DeclareMathOperator*\uplim{\overline{lim}}
\DeclareMathOperator\rng{rng}
\DeclareMathOperator\Sym{Sym}
\DeclareMathOperator\Arg{Arg}
\DeclareMathOperator\Log{Log}
\DeclareMathOperator\dom{dom}
\DeclareMathOperator\supp{Supp}
\DeclareMathOperator\var{Var}
\DeclareMathOperator\cov{Cov}

% commands
%\renewcommand\qedsymbol{\textbf{מש''ל}}
%\renewcommand\qedsymbol{\fbox{\emoji{lizard}}}
\newcommand{\Aa}[0]{\mathcal{A}}
\newcommand{\Bb}[0]{\mathcal{B}}
\newcommand{\CC}[0]{\mathbb{C}}
\newcommand{\Cc}[0]{\mathcal{C}}
\newcommand{\EE}[0]{\mathbb{E}}
\newcommand{\FF}[0]{\mathbb{F}}
\newcommand{\Ff}[0]{\mathcal{F}}
\newcommand{\Ii}[0]{\mathcal{I}}
\newcommand{\Gg}[0]{\mathcal{G}}
\newcommand{\Ll}[0]{\mathcal{L}}
\newcommand{\Mm}[0]{\mathcal{M}}
\newcommand{\NN}[0]{\mathbb{N}}
\newcommand{\Nn}[0]{\mathcal{N}}
\newcommand{\PP}[0]{\mathbb{P}}
\newcommand{\Pp}[0]{\mathcal{P}}
\newcommand{\QQ}[0]{\mathbb{Q}}
\newcommand{\RR}[0]{\mathbb{R}}
\newcommand{\Rr}[0]{\mathcal{R}}
\newcommand{\Ss}[0]{\mathcal{S}}
\newcommand{\TT}[0]{\mathbb{T}}
\newcommand{\Uu}[0]{\mathcal{U}}
\newcommand{\Vv}[0]{\mathcal{V}}
\newcommand{\Ww}[0]{\mathcal{W}}
\newcommand{\ZZ}[0]{\mathbb{Z}}
\newcommand{\acts}[0]{\circlearrowright}
\newcommand{\explain}[2] {
	\begin{flalign*}
		 && \text{#2} && \text{#1}
	\end{flalign*}
}
\newcommand{\maketitleprint}[0]{ \begin{center}
	%\begin{tikzpicture}[scale=3]
	%	\duck[graduate=gray!20!black, tassel=red!70!black]
	%\end{tikzpicture}	
	\includegraphics[width=6cm]{cover}
\end{center}
}

% theorem commands
\newtheoremstyle{c_remark}
	{}	% Space above
	{}	% Space below
	{}% Body font
	{}	% Indent amount
	{\bfseries}	% Theorem head font
	{}	% Punctuation after theorem head
	{.5em}	% Space after theorem head
	{\thmname{#1}\thmnumber{ #2}\thmnote{ \normalfont{\text{(#3)}}}}	% head content
\newtheoremstyle{c_definition}
	{3pt}	% Space above
	{3pt}	% Space below
	{}% Body font
	{}	% Indent amount
	{\bfseries}	% Theorem head font
	{}	% Punctuation after theorem head
	{.5em}	% Space after theorem head
	{\thmname{#1}\thmnumber{ #2}\thmnote{ \normalfont{\text{(#3)}}}}	% head content
\newtheoremstyle{c_plain}
	{3pt}	% Space above
	{3pt}	% Space below
	{\itshape}% Body font
	{}	% Indent amount
	{\bfseries}	% Theorem head font
	{}	% Punctuation after theorem head
	{.5em}	% Space after theorem head
	{\thmname{#1}\thmnumber{ #2}\thmnote{ \text{(#3)}}}	% head content

\ifcsname c@english\endcsname
	\theoremstyle{plain}
	\newtheorem{theorem}{Theorem}[section]
	\newtheorem{lemma}[theorem]{Lemma}
	\newtheorem{proposition}[theorem]{Proposition}
	\newtheorem*{proposition*}{Proposition}
	%\newtheorem{corollary}[theorem]{אין חלופה עברית}

	\theoremstyle{definition}
	\newtheorem{definition}[theorem]{Definition}
	\newtheorem*{definition*}{Definition}
	\newtheorem{example}{Example}[section]
	\newtheorem{exercise}{Exercise}[section]

	\theoremstyle{remark}
	\newtheorem*{remark}{Remark}
	\newtheorem*{solution}{Solution}
	\newtheorem{conclusion}[theorem]{Conclusion}
	\newtheorem{notation}[theorem]{Notation}
\else
	\theoremstyle{c_plain}
	\newtheorem{theorem}{משפט}[section]
	\newtheorem{lemma}[theorem]{למה}
	\newtheorem{proposition}[theorem]{טענה}
	\newtheorem*{proposition*}{טענה}
	%\newtheorem{corollary}[theorem]{אין חלופה עברית}

	\theoremstyle{c_definition}
	\newtheorem{definition}[theorem]{הגדרה}
	\newtheorem*{definition*}{הגדרה}
	\newtheorem{example}{דוגמה}[section]
	\newtheorem{exercise}{תרגיל}[section]

	\theoremstyle{c_remark}
	\newtheorem*{remark}{הערה}
	\newtheorem*{solution}{פתרון}
	\newtheorem{conclusion}[theorem]{מסקנה}
	\newtheorem{notation}[theorem]{סימון}
\fi

% Questions related commands
\newcounter{question}
\setcounter{question}{1}
\newcounter{sub_question}
\setcounter{sub_question}{1}

\ifcsname c@english\endcsname
	\newcommand{\question}[1][0]{
		\ifthenelse{#1 = 0}{}{\setcounter{question}{#1}}
		\section{Question \arabic{question}}
		\addtocounter{question}{1}
		\setcounter{sub_question}{1}
	}

	\newcommand{\subquestion}[1][0]{
		\ifthenelse{#1 = 0}{}{\setcounter{sub_question}{#1}}
		\subsection{Part \alph{sub_question}}
		\addtocounter{sub_question}{1}
	}
\else
	\newcommand{\question}[1][0]{
		\ifthenelse{#1 = 0}{}{\setcounter{question}{#1}}
		\section{שאלה \arabic{question}}
		\addtocounter{question}{1}
		\setcounter{sub_question}{1}
	}

	\newcommand{\subquestion}[1][0]{
		\ifthenelse{#1 = 0}{}{\setcounter{sub_question}{#1}}
		\subsection{סעיף \localecounter{letters.gershayim}{sub_question}}
		\addtocounter{sub_question}{1}
	}
\fi

% import lua and start of document
\directlua{common = require ('../common')}

\GetEnv{AUTHOR}

% headers
\author{\AUTHOR}
\date\today

\title{פתרון מטלה 09 --- אנליזה על יריעות, 80426}
% chktex-file 17
% chktex-file 9

\DeclareMathOperator{\vol}{vol}

\begin{document}
\maketitle
\maketitleprint[blue]

\question{}
תהי $M^k \subseteq \RR^n$ יריעה עם שפה.
נסמן $M^\circ = M \setminus \partial M$.

\subquestion{}
נראה ש־$M^\circ$ היא צפופה ב־$M$.
\begin{proof}
	$M^\circ$ צפופה ב־$M$ אם ורק אם כל נקודה ב־$M$ היא ערך גבולי של סדרה ב־$M \setminus \partial M$.
	עבור נקודות ב־$M$ הטענה נכונה באופן טריוויאלי, לכן מספיק שנמצא סדרה ${\{ x_n \}}_{n = 1}^\infty \subseteq M^\circ$ המתכנסת ל־$x \in \partial M$ לכל $x$ כזה.
	נניח ש־$x \in \partial M$ ונניח ש־$\alpha : U \to M$ פרמטריזציה מקומית של $x$, כלומר $p \in U \subseteq \HH^k$ כאשר $p = \alpha^{-1}(x)$.
	פרמטריזציה משמרת שפה, כלומר $y \in \partial U \iff \alpha(y) \in \partial M$, אבל $U$ פתוחה ולכן קיים $\epsilon > 0$ כך ש־$B(p, \epsilon) \cap \HH^k \subseteq U$.
	נבחר נקודה $x_1 \in B(p, \epsilon) \cap U \cap {(\HH^k)}^\circ$ וכן נקודה $x_n \in B(p, \frac{\epsilon}{n}) \cap U \cap {(\HH^k)}^\circ$ לכל $n \in \NN$.
	נקבל סדרה כך ש־$x_n \to x$ וכן $x_n \in M^\circ$ לכל $n \in \NN$, כפי שרצינו.
\end{proof}

\subquestion{}
נראה שאם $k = n$ אז $M^\circ$ היא קבוצה פתוחה ב־$\RR^n$.
\begin{proof}
	ראינו בהרצאה ש־$M^\circ$ היא יריעה ללא שפה ממימד $k = n$, וכן ראינו שיריעות $M^n \subseteq \RR^n$ הן קבוצות פתוחות ב־$\RR^n$.

	הטענה נובעת באופן ישיר מקיום פרמטריזציה מקומית.
	אם $x \in M^\circ$ אז קיימת פרמטריזציה $\alpha : U \to M^\circ$ כך ש־$U \subseteq \RR^n$ פתוחה, ו־$\alpha$ הומיאומורפיזם ולכן בפרט פתוחה ו־$\alpha(U) \subseteq M$ פתוחה במובן המטרי.
	נסיים את ההוכחה עם סגירות פתוחות לאיחוד ואיחוד כלל הסביבות הפתוחות לכל נקודה ב־$M^\circ$.
\end{proof}

\subquestion{}
נראה שאם $U \subseteq M$ קבוצה פתוחה ב־$\RR^n$ אז $U \subseteq M^\circ$.
\begin{proof}
	נתון כי $U \subseteq M$ ו־$\deg U = n$ ולכן $\deg M = n$.
	לכל נקודה $x \in U$ יש פרמטריזציה מקומית וכמו בסעיף הקודם היא פתוחה ולכן מעידה ש־$x \in M^\circ$ כקבוצה פתוחה (הנחה שמותר לעשות בשל תוצאת הסעיף הקודם).
\end{proof}

\subquestion{}
נסיק שאם $M$ תת־קבוצה סגורה של $\RR^n$ ו־$\dim M = n$ אז ההגדרות של שפה ופנים של יריעות ושל קבוצות מזדהות.
\begin{proof}
	נוכל להסיק ישירות מהסעיף הקודם שפנים של יריעה וקבוצה מזדהים כאשר $\dim M = n$.
	ביתר פירוט, אם $U \subseteq M$ פתוחה אז בפרט $U \subseteq M^\circ$ ולכן $M^\circ$ תת־קבוצה פתוחה מקסימלית של $M$, כלומר היא עומדת בהגדרה המטרית המדויקת של פנים.
	נקבל על־ידי חיסור קבוצות והעובדה ש־$M = \overline{M}$ ש־$\partial M = \overline{M} \setminus M^\circ$, כלומר גם השפה מזדהה עם ההגדרה המטרית.
\end{proof}

\question{}
תהי $M^k \subseteq \mathbb{R}^n$ יריעה עם שפה,
נניח ש־$q \in \partial M$ ותהי פרמטריזציה $\alpha : V \to W$ כך ש־$V \subseteq \HH^k$ פתוחה ו־$q \in W \subseteq M$.
נניח גם ש־$\tilde{\alpha} : \tilde{V} \to \tilde{W}$ הרחבה חלקה שלה, כלומר $\tilde{V} \subseteq \RR^k$ ו־$\dim \tilde{W} = k$. \\
נניח ש־$y : [0, \varepsilon) \to \tilde{V}$ היא מסילה חלקה כך ש־$y(0) = y_0 = \alpha^{-1}(q)$, וש־$d \tilde{\alpha}_{y_0}(y'(0)) = X(q)$ מקיימת $\langle X(q), \nu(q) \rangle < 0$,
כאשר $\nu$ הנורמל החיצוני של $M$. \\
נראה שקיים $\delta > 0$ כך ש־$y(t) \in \HH^k$ עבור $t \in [0, \delta]$,
ובפרט ש־$\tilde{\alpha}(y(t)) = \alpha(y(t)) \in M$.
\begin{proof}
	אנו נראה שהקורדינטה האחרונה של $y'(0)$ חיובית.
	נפעיל תהליך גרם־שמידט על $X(q)$ ונקבל,
	\[
		X(q)
		= \langle X(q), \nu(q) \rangle \nu(q) + Y(q)
	\]
	כאשר $Y(q) \in T_q \partial M$ מאורתוגונליות הנורמל החיצוני והמשיק לשפה.
	בהתאם הקורדינטה האחרונה של $d Y(q)$ אפס כווקטור בנגזרת השפה, $\langle X(q), \nu(q) \rangle$ ערך שלילי ו־$\nu(q)$ שלילי בקורדינטה האחרונה, ולכן $X(q)$ חיובי בקורדינטה האחרונה.

	עתה נסיק כי קיים $\delta > 0$ כך ש־$y(t) \in \HH^k$, זאת שכן הוא נע מעלה בקורדינטה האחרונה ומוגדר בכל אלו שאינן האחרונות.
	בפרט $\tilde{\alpha}(y(t)) = \alpha(y(t))$, שכן מהנתון ההרחבה היא עבור ערכים שליליים.
\end{proof}

\question{}
\subquestion{}
תהי $M^k \subseteq \RR^n$ יריעה קומפקטית עם שפה.
נניח ש־$f, g : M \to \RR$ רציפות וש־$c \in \RR$.
נראה שמתקיים,
\[
	\int_M (\alpha f + g)\ d\vol_k
	\alpha \int_M f\ d\vol_k + \int_M g\ d\vol_k
\]
\begin{proof}
	נניח ש־${\{ \alpha_i : U_i \to W_i \}}_{i = 1}^N$ פרמטריזציות מקומיות המכסות את $M$, כלומר $M = \bigcup_{i = 1}^N W_i$.
	נניח גם ש־${\{ \varphi_i : M \to W_i \}}_{i = 1}^N$ חלוקת יחידה עבור $\{ W_i \}$.
	אז מהגדרת האינטגרל,
	\[
		\int_M (c f + g)\ d\vol_k
		= \sum_{i = 1}^N \int_{U_i} \varphi_i(\alpha_i(x)) (c f + g)(\alpha_i(x)) V(D \alpha_i |_x)\ dx
	\]
	נבחין כי זהו אינטגרל רב משתני רגיל ולכן לינארי מתכונות האינטגרל, כלומר,
	\begin{multline*}
		\sum_{i = 1}^N \int_{U_i} \varphi_i(\alpha_i(x)) (c f + g)(\alpha_i(x)) V(D \alpha_i |_x)\ dx \\
		= \sum_{i = 1}^N \left( c \int_{U_i} \varphi_i(\alpha_i(x)) f(\alpha_i(x)) V(D \alpha_i |_x)\ dx + \int_{U_i} \varphi_i(\alpha_i(x)) g(\alpha_i(x)) V(D \alpha_i |_x)\ dx \right)
	\end{multline*}
	אבל גם סכומים סופיים הם לינאריים ונסיק,
	\begin{multline*}
		\sum_{i = 1}^N \left( c \int_{U_i} \varphi_i(\alpha_i(x)) f(\alpha_i(x)) V(D \alpha_i |_x)\ dx + \int_{U_i} \varphi_i(\alpha_i(x)) g(\alpha_i(x)) V(D \alpha_i |_x)\ dx \right) \\
		= \left( c \sum_{i = 1}^N \int_{U_i} \varphi_i(\alpha_i(x)) f(\alpha_i(x)) V(D \alpha_i |_x)\ dx \right) + \sum_{i = 1}^N \left( \int_{U_i} \varphi_i(\alpha_i(x)) g(\alpha_i(x)) V(D \alpha_i |_x)\ dx \right)
	\end{multline*}
	והביטוי האחרון איננו אלא,
	\[
		\alpha \int_M f\ d\vol_k + \int_M g\ d\vol_k
	\]
	כפי שרצינו.
\end{proof}

\subquestion{}
משפט פוביני,
נניח ש־$M^m$ ו־$N^n$ יריעות, ונניח כי $f, g : M \times N \to \RR$ פונקציה רציפה,
אז,
\[
	\int_{M \times N} f\ d\vol_{m + n}
	= \int_{M} \left( \int_N f\ d\vol_n \right)\ d\vol_m
	= \int_{N} \left( \int_M f\ d\vol_m \right)\ d\vol_n
\]
\begin{proof}
	נניח ש־${\{ \alpha_i : U_i \to W_i \}}_{i = 1}^k$ פרמטריזציות מקומיות המכסות את $M \times N$, כלומר $M \times N = \bigcup_{i = 1}^k W_i$. \\
	נניח גם ש־${\{ \varphi_i : M \times N \to W_i \}}_{i = 1}^k$ חלוקת יחידה עבור $\{ W_i \}$.
	נניח כי $U_i = U_i^M \times U_i^N$ עבור $U_i^M \subseteq \HH^m, U_i^N \subseteq \HH^n$ לכל $i \le k$. \\
	נסמן פירוק דומה ל־$W_i$, וכן עבור $\alpha_i, \varphi_i$ לכל $i$. \\
	מהגדרה,
	\begin{align*}
		\int_{M \times N} f\ d\vol_{m + n}
		& = \sum_{i = 1}^k \int_{U_i} \varphi_i(\alpha_i(x)) f(\alpha_i(x)) V(D \alpha_i |_x)\ dx \\
		& = \sum_{i = 1}^k \int_{U_i^M \times U_i^N} \varphi_i(\alpha_i(x, y)) f(\alpha_i(x, y)) V(D \alpha_i |_{(x, y)})\ d(x, y)
	\end{align*}
	כאשר $x \in \RR^m, y \in \RR^n$ בביטוי האחרון.
	זהו סכום סופי של אינטגרלים רב־משתניים ולכן משפט פוביני חל בהם, נראה את אחת הגרסות כדי לחסוך בנייר,
	\begin{align*}
		& \sum_{i = 1}^k \int_{U_i^M \times U_i^N} \varphi_i(\alpha_i(x, y)) f(\alpha_i(x, y)) V(D \alpha_i |_{(x, y)})\ dx\ dy \\
		\overset{(1)}{=}  & \sum_{i = 1}^k \int_{U_i^M} \left( \int_{U_i^N} \varphi_i(\alpha_i(x, y)) f(\alpha_i(x, y)) V(D \alpha_i |_{(x, y)})\ dy \right)\ dx \\
		\overset{(2)}{=}  & \sum_{i = 1}^k \int_{U_i^M} \left( \sum_{j = 1}^k \int_{U_i^N} \varphi_i^M(\alpha_i^M(x)) \varphi_j^N(\alpha_j^N(y)) f(\alpha_i(x, y)) V(D \alpha_i |_{(x, y)})\ dy \right)\ dx \\
		\overset{(3)}{=}  & \sum_{i = 1}^k \int_{U_i^M} \left( \varphi_i^M(\alpha_i^M(x)) \sum_{j = 1}^k \int_{U_i^N} \varphi_j^N(\alpha_j^N(y)) f(\alpha_i(x, y)) V(D \alpha_i^N |_{y})\ dy \right)V(D \alpha_i^M |_{(x)})\ dx \\
		= & \sum_{i = 1}^k \int_{U_i^M} \left( \varphi_i^M(\alpha_i^M(x)) \int_{N} f(x, y)\ d\vol_n \right) V(D \alpha_i^M |_{(x)})\ dx \\
		= & \int_{M} \left( \int_{N} f\ d\vol_n \right)\ d\vol_m
	\end{align*}
	כאשר,
	\begin{enumerate}
		\item משפט פוביני לאינטגרלים רב־משתניים
		\item הוספת סכום ומציין יחידה
		\item כפליות דטרמיננטת בלוקים
	\end{enumerate}
\end{proof}

\question{}
תהי $M^k$ יריעה קומפקטית ונניח ש־$M_1^k, M_2^k$ שתי יריעות קומפקטיות עם שפה המקיימות,
\[
	M_1 \cup M_2 = M,
	\quad
	M_1 \cap M_2
	= \partial M_1
	= \partial M_2
	= \Gamma
\]
נניח ש־$f : M \to \RR$ רציפה,
אנו נראה כי,
\[
	\int_M f\ d\vol_k
	= \int_{M_1} f\ d\vol_k
	+ \int_{M_2} f\ d\vol_k
\]
\begin{proof}
	נניח ש־${\{ \alpha_i : U_i \to W_i \}}_{i = 1}^l$ פרמטריזציות מקומיות של $M$ המכסות אותה.
	נניח גם כי כל פרמטריזציה מושרית מנקודה ב־$\Gamma$ או זרה ל־$\Gamma$.
	מותר לנו להניח כך שכן נוכל לקחת כיסוי פתוח מפרמטריזציות של $\Gamma$ (היא קומפקטית) והרחבת הכיסוי לכל $M$ ושימוש בחיתוכים סופיים.
	נניח ש־${\{ \varphi_i : M \to W_i \}}_{i = 1}^l$ כיסוי יחידה הכפוף ל־$\{ W_i \}$.
	מההנחה הנוספת נוכל גם להניח בלי הגבלת הכלליות כי קיימים $l_1 < l_2 < l$ כך ש־$W_i \subseteq M_1$ עבור $1 \le i \le l_1$,
	$W_i \subseteq M_2$ עבור $l_1 < i \le l_2$ ולכל $l_2 < i \le l$ ש־$\alpha_i$ מושרית מנקודה ב־$\Gamma$.

	ממשפט מהתרגול לכל $l_2 < i \le l$, הפרמטריזציה $\alpha_i$ מקיימת,
	\[
		\alpha_i(\HH^k \cap U_i) \subseteq M_1,
		\qquad
		\alpha_i((-\HH^k) \cap U_i) \subseteq M_2
	\]
	ונסמן את שתי הפרמטריזציות המקומיות המושרות הללו $\alpha_i^1, \alpha_i^2$ בהתאמה. \\
	נגדיר גם פירוק של $\varphi_i$ ל־$\varphi_i^1, \varphi_i^2$ על־ידי מכפלה במציין $1_{M_1}$ ובמציין $1_{M_2}$ בהתאמה.
	עתה אנו מוכנים ויכולים לבחון את האינטגרל,
	\begin{align*}
		& \int_M f\ d\vol_k \\
		= & \sum_{i = 1}^l \int_{U_i} \varphi_i(\alpha_i(x)) f(\alpha_i(x)) V(D \alpha_i |_{x})\ dx \\
		= & \sum_{i = 1}^{l_1} \int_{U_i} \varphi_i(\alpha_i(x)) f(\alpha_i(x)) V(D \alpha_i |_{x})\ dx \\
		  & + \sum_{i = l_1 + 1}^{l_2} \int_{U_i} \varphi_i(\alpha_i(x)) f(\alpha_i(x)) V(D \alpha_i |_{x})\ dx \\
		  & + \sum_{i = l_2 + 1}^{l} \int_{U_i} \varphi_i^1(\alpha_i(x)) f(\alpha_i(x)) V(D \alpha_i |_{x})\ dx \\
		  & + \sum_{i = l_2 + 1}^{l} \int_{U_i} \varphi_i^2(\alpha_i(x)) f(\alpha_i(x)) V(D \alpha_i |_{x})\ dx \\
		= & \sum_{1 \le i \le l_1, l_2 < i \le l} \int_{U_i} \varphi_i^1(\alpha_i(x)) f(\alpha_i(x)) V(D \alpha_i |_{x})\ dx \\
		  & + \sum_{l_1 < i \le l} \int_{U_i} \varphi_i^2(\alpha_i(x)) f(\alpha_i(x)) V(D \alpha_i |_{x})\ dx \\
		= & \sum_{1 \le i \le l_1, l_2 < i \le l} \int_{U_i^1} \varphi_i^1(\alpha_i^1(x)) f(\alpha_i^1(x)) V(D \alpha_i^1 |_{x})\ dx \\
		  & + \sum_{l_1 < i \le l} \int_{U_i^2} \varphi_i^2(\alpha_i^2(x)) f(\alpha_i^2(x)) V(D \alpha_i^2 |_{x})\ dx \\
		= & \int_{M_1} f\ d\vol_k + \int_{M_2} f\ d\vol_k
	.\end{align*}
\end{proof}

\end{document}
